%======================================================================
\section{Notación}
%======================================================================
\begin{table}[H]\centering\small
\begin{tabularx}{\linewidth}{@{}l Y@{}}\toprule
Símbolo & Significado\\\midrule
\multicolumn{2}{@{}l}{\emph{Reactor}}\\
$N$, $C_i$, $C^e_i$ & Población de neutrones; precursores dentro y fuera del núcleo (6 grupos)\\
$\rho$, $\rho_0$, $\beta$, $\Lambda$ & Reactividad; reactividad estacionaria; fracción de neutrones retardados; tiempo de generación\\
$\lambda_i$, $\lambda_c$, $\lambda_e$ & Constantes de decaimiento de los precursores; tasas de salida y de regreso al núcleo\\
$\alpha_f$, $\alpha_c$ & Coeficientes de realimentación de temperatura del combustible y del refrigerante\\
$T_f$, $T_w$, $T_c$ & Temperaturas medias del combustible, la pared y el plomo\\
$M_c$, $\Delta M$ & Inventario nominal de plomo; inventario ganado (negativo: perdido)\\
$\mdot_c$, $\mdot_{\rm nom}$ & Caudal primario actual y nominal (5550.5~kg/s)\\
$h_{fw}$, $h_{wc}$, $A_{fw}$, $A_{wc}$ & Coeficientes y áreas de transferencia de calor combustible--pared y pared--plomo\\
$p_4$ & Polinomio de corrección de la PCT (de la media al pico)\\
\multicolumn{2}{@{}l}{\emph{Control}}\\
$x$, $\delta_k$ & Estado del ROM (17 componentes); desviación respecto de una trayectoria nominal\\
$u_k=\mdot_{{\rm ACS},k}$ & Caudal del ACS en el paso $k$ (la decisión), en kg/s\\
$\hat d_k=\mdot_{{\rm leak},k}$ & Pronóstico de la fuga\\
$\hat y_k$, $y_{\rm ref}$, $U_y$ & PCT estimada, referencia y límite de seguridad\\
$K$, $\Delta t$ ($h$), $t_q$ & Horizonte (26 pasos), paso de tiempo (0.5~s), instante de la consulta\\
$\lambda_m$, $\lambda_{\rm pct}$, $\Delta\mdot$ & Pesos de las penalizaciones de suavidad y de seguridad; salto de caudal admitido\\
$\lambda^{\rm A}_j$, $\mu^{\rm A}_j$ & Tasas de falla y de reparación de los componentes del ACS\\
\multicolumn{2}{@{}l}{\emph{Optimización}}\\
$J$, $J^\star$ & Costo P0 y su óptimo global\\
$UB$, $LB$ & Cotas superior e inferior de $J^\star$\\
$A$, $B$, $E$, $D$ & Jacobianos de la dinámica respecto del estado, el control y la fuga; derivada de $p_4$\\
$G$, $G(s)$ & Matriz condensada de sensibilidad de la salida al control (26$\times$26); función de transferencia\\
$[A](Z)$ & Jacobiano evaluado con aritmética de intervalos sobre la caja $Z$\\
$\varepsilon_k$ & Cota del residuo no lineal en el paso $k$\\
N1, N2, N3 & Formulaciones LTI, LTV iterada y no lineal\\
BE, \code{ode1be} & Euler implícito convergido; Euler implícito con una iteración de Newton (Simulink)\\\bottomrule
\end{tabularx}
\end{table}

%======================================================================
\section{Glosario}
%======================================================================
\begin{description}[style=nextline, leftmargin=0pt, itemsep=3pt]
  \item[ACS] Sistema auxiliar de enfriamiento: repone plomo en el circuito primario con una
        bomba MHD.
  \item[A-estable] Propiedad de un integrador numérico que no amplifica los modos que decaen,
        cualquiera sea el paso.
  \item[BiLSTM] Red neuronal recurrente bidireccional con celdas de memoria de corto y largo
        plazo; aquí pronostica la fuga.
  \item[Brecha de optimalidad] Diferencia relativa entre una cota superior (el costo de una
        solución) y una cota inferior demostrada del óptimo global.
  \item[Cinética puntual] Modelo de la población de neutrones que supone una forma espacial
        fija del flujo.
  \item[Condiciones KKT] Condiciones de primer orden que cumple todo mínimo local regular de un
        problema con restricciones.
  \item[Convexo] Problema cuyo conjunto factible y costo son convexos; en él, todo mínimo local
        es global.
  \item[DET] Árbol dinámico de eventos: simulación que se ramifica cuando un componente falla o
        se repara.
  \item[Diferenciación automática] Cálculo exacto de derivadas propagando la regla de la cadena
        por las operaciones de un programa.
  \item[Disparo simple] Formulación de control óptimo en la que sólo los controles son
        variables y el costo se obtiene simulando.
  \item[HRM] Modelo de alta resolución (12 nodos) del SMDFR; hace de planta.
  \item[LOCA] Accidente con pérdida de refrigerante.
  \item[Metaheurística] Método de búsqueda general que sólo evalúa el costo (LJ, GA, DE, PSO).
  \item[MC-dropout] Estimación de la incertidumbre de una red neuronal repitiendo la inferencia
        con dropout activo.
  \item[MPC] Control predictivo basado en modelo: optimiza un horizonte futuro, aplica la primera
        acción y repite.
  \item[NLP / QP] Programa no lineal / programa cuadrático.
  \item[PCT] Temperatura pico de la vaina (aquí, de la pared de los tubos): el indicador de
        seguridad.
  \item[Precursores] Fragmentos de fisión que emiten neutrones retardados al decaer.
  \item[PSM / DSM] Margen probabilístico / determinista de seguridad.
  \item[Reactividad] Medida de cuánto se aparta el reactor de la criticidad,
        $\rho=(k-1)/k$.
  \item[ROM] Modelo reducido (un nodo por material) usado por el controlador para predecir.
  \item[SAA] Aproximación de una esperanza por el promedio sobre escenarios muestreados.
  \item[Transcripción directa] Formulación monolítica: estados y controles son variables, y la
        dinámica discretizada es un conjunto de restricciones.
\end{description}

%======================================================================
\section{Reproducibilidad}
%======================================================================
\begin{itemize}\sloppy
  \item \textbf{Réplica y NLP}: \code{julia/SMDFRROM.jl}, \code{validate\_replica.jl},
        \code{f2\_monolithic\_nlp.jl}.
  \item \textbf{QP, certificados y forma condensada}: \code{julia/f3\_certificates\_condensed.jl}.
  \item \textbf{Certificación}: \code{julia/ReducedP0.jl}, \code{f3c\_optimality.jl}
        (pasos 1--3), \code{f3c\_lower\_bound.jl} (4a), \code{f3c\_rigorous.jl} (4b; centro
        con \code{F3C\_CENTER=mid}), \code{f3c\_compare\_decisions.jl}.
  \item \textbf{Riesgo y actuador}: \code{python/f4\_leak\_ensemble.py},
        \code{f4\_forecast\_calibration.py}, \code{julia/f4\_risk\_decisions.jl},
        \code{f4\_actuator\_sweep.jl}.
  \item \textbf{Análisis didácticos} (valores propios, respuesta en frecuencia, curva de la
        bomba, autoridad): \code{julia/didactic\_analysis.jl} $\Rightarrow$
        \code{results/didactic/}.
  \item \textbf{Resultados}: \code{results/f1\_rom\_replica/}, \code{f3\_certificates/},
        \code{f3c\_optimality/}, \code{f4\_risk/}, \code{didactic/}.
  \item \textbf{Figuras}: \code{python/make\_status\_figures.py} (resultados) y
        \code{python/make\_didactic\_figures.py} (capítulos introductorios). Las figuras de
        rigidez y de brecha de optimalidad en una dimensión son ejemplos ilustrativos, no
        resultados del modelo.
  \item \textbf{Documento}: \code{tectonic -X compile estado\_mpc\_monolitico.tex} desde
        \code{docs/}; los capítulos están en \code{docs/estado/}.
\end{itemize}
